
These datasets will be used to develop and evaluate a new generation of
machine-learning-based closure models. In particular, Physics-Informed Neural
Networks (PINNs) approaches will be employed to integrate information from
experiments, system-scale calculations, and high-fidelity simulations into
unified predictive frameworks. The resulting closures will be benchmarked
against both conventional RANS approaches and experimental measurements,
providing a quantitative assessment of their accuracy, robustness, and
generalizability across operating conditions. The ultimate goal is to create
data-informed closures capable of significantly improving the predictive
capability of engineering-scale reactor simulations while maintaining
computational efficiency.

The proposed work will also advance the development of NekROM
\cite{nekrom,kaneko20a,kaneko22a,tsai24a}, a reduced-order
modeling framework built upon high-fidelity Nek5000/RS simulations.
NekROM leverages techniques such as Proper Orthogonal Decomposition (POD),
operator inference, and machine-learning-assisted model reduction to construct
low-dimensional representations of complex transient flows. By extracting the
dominant flow structures and dynamics from DNS and LES datasets, NekROM enables
the prediction of transient behavior at a fraction of the computational cost of
the original simulations. The expanded transient database generated through
this project provides a unique opportunity to further develop, train, and
validate NekROM for reactor-relevant applications involving mixed convection,
thermal stratification, and natural circulation.

More broadly, the project seeks to establish a systematic data-driven workflow
for translating information from high-fidelity simulations into practical
engineering tools. Emerging technologies such as PINNs,
machine-learning-assisted closure development, and advanced reduced-order
models will be integrated within a common multiscale framework. Together, these
capabilities will create a robust pathway for incorporating high-fidelity
transient physics into reactor design and safety-analysis tools, significantly
advancing the state of the art in multiscale thermal-hydraulic modeling and
accelerating the deployment of predictive simulation capabilities within the
DOE NEAMS ecosystem.


@article{kaneko22a,
author = "Kaneko, K. and P. Fischer",
title = "Augmented Reduced Order Models for Turbulence",
journal = "Frontiers in Physics",
year = "2022"
}
@article{tsai24a,
author = "Tsai, P.H. and P. Fischer and T. Iliescu",
title = "A Time-Relaxation Reduced Order Model for Turbulent Channel Flow",
journal = jcp,
volume = "521",
pages="113563",
year = 2025}
@article{kaneko20a,
author = "Kaneko, K and P.-H. Tsai and P. Fischer",
title = "Towards Model Order Reduction for Fluid-Thermal Analysis",
journal = "Nuc. Eng. and Des.",
volume = "370",
pages = "110866",
publisher = "Elsevier BV",
year = "2020"
}
@MISC{nekrom,
key   = {NekROM},
author  = {Kaneko, K. and P.H. Tsai and N. Christensen and P. Fischer and E. Merzari},
title = "NekROM: Reduced-Order Models for {N}ek5000/{RS}",
howpublished = "{\url{github.com/Nek5000/NekROM}}",
year = "2020"
}
